%%%%%%%%%%%%%%%%%%%%%%%%%%%%%%%%%%%%%%%%%%%%%%%%%%%%%%%%%%%%%%%%%%%%%%%%%%%%%%%%
% AMS Beamer series / Bologna FC / Template
% Andrea Omicini
% Alma Mater Studiorum - Università di Bologna
% mailto:andrea.omicini@unibo.it
%
% This template is released into the public domain under CC0 1.0 Universal:
% use it, adapt it and build on it freely, with no obligation to credit.
% See the LICENSE file, or https://creativecommons.org/publicdomain/zero/1.0/
%
% The AMSBolognaFC style it uses is covered by its own licence (LPPL 1.3c):
% https://github.com/andreaomicini/beamer-AMSBolognaFC
%%%%%%%%%%%%%%%%%%%%%%%%%%%%%%%%%%%%%%%%%%%%%%%%%%%%%%%%%%%%%%%%%%%%%%%%%%%%%%%%
\documentclass[presentation,cesena,apice]{beamer}\mode<presentation>{\usetheme{AMSBolognaFC}}
%\documentclass[presentation]{beamer}\mode<presentation>{\usetheme{AMSBolognaFC}}
%\documentclass[handout]{beamer}\mode<handout>{\usetheme{AMSBolognaFC}}
%%%%%%%%%%%%%%%%%%%%%%%%%%%%%%%%%%%%%%%%%%%%%%%%%%%%%%%%%%%%%%%%%%%%%%%%%%%%%%%%
\def\BibTeX{{\rm B\kern-.05em{\sc i\kern-.025em b}\kern-.08em
    T\kern-.1667em\lower.7ex\hbox{E}\kern-.125emX}}
\renewcommand{\thefootnote}{\fnsymbol{footnote}}
% a colour swatch, framed so the pale ones still read against white
\newcommand{\swatch}[1]{\fcolorbox{bolognafcsilver!70!black}{#1}{\rule{0pt}{1.3ex}\rule{2.2ex}{0pt}}}
% version -- Major.Minor[.Patch]; the release time-stamp is appended by the CI
\newcommand{\templatemajor}{1}
\newcommand{\templateminor}{10}
% \templatepatch is optional: comment the next line out to release as Major.Minor
\newcommand{\templatepatch}{3}
\ifdefined\templatepatch{}
  \newcommand{\templateversion}{\templatemajor.\templateminor.\templatepatch}
\else
  \newcommand{\templateversion}{\templatemajor.\templateminor}
\fi
%%%%%%%%%%%%%%%%%%%%%%%%%%%%%%%%%%%%%%%%%%%%%%%%%%%%%%%%%%%%%%%%%%%%%%%%%%%%%%%%
\title[AMS Bologna FC Beamer Style]
{AMS Bologna FC Beamer Style}
%
\subtitle[AMS Series Templates]
{AMS Series Templates}
%
\author[\sspeaker{Omicini} \and Else]
{\speaker{Andrea Omicini} \and Someone Else}
%
\institute[DISI, Univ. Bologna]
{\url{mailto:andrea.omicini@unibo.it} \qquad \url{mailto:someone.else@unibo.it}\\
Dipartimento di Informatica \textendash{} Scienza e Ingegneria (DISI)\\
\textsc{Alma Mater Studiorum} \textendash{} Università di Bologna}
%
\date[v.\ \templateversion]{version \templateversion}
%
%%%%%%%%%%%%%%%%%%%%%%%%%%%%%%%%%%%%%%%%%%%%%%%%%%%%%%%%%%%%%%%%%%%%%%%%%%%%%%%%
\begin{document}
%%%%%%%%%%%%%%%%%%%%%%%%%%%%%%%%%%%%%%%%%%%%%%%%%%%%%%%%%%%%%%%%%%%%%%%%%%%%%%%%

%/////////
\frame{\titlepage}
%/////////

%%===============================================================================
%\section*{Outline}
%%===============================================================================
%
%%/////////
%\frame[c]{\tableofcontents[hideallsubsections]}
%%/////////

%===============================================================================
\section{Premises}
%===============================================================================

%/////////
\begin{frame}[c]{Academic Presentations}
%
\begin{itemize}
	\item technical presentations are common in the academic environment
	\item \alert{\LaTeX{}}\footnote{\uurl{https://www.latex-project.org}} is typically used for scientific literature
	\begin{itemize}
		\item for better formulas, bibliography, typography, \ldots
	\end{itemize}
	\item so, \LaTeX{} is the most typical \emph{toolbox} for scientific presentations
	\begin{itemize}
		\item which easily translates into teaching presentations, too
	\end{itemize}
\end{itemize}
%
\end{frame}
%/////////

%/////////
\begin{frame}[c]{Beamer}
%
\begin{itemize}
	\item \alert{Beamer} is a \LaTeX{} document class for creating presentation \emph{slides}, with a wide range of \emph{templates} and a set of features for making \emph{slideshow effects}
	\begin{description}
		\item[repo] \uurl{https://github.com/josephwright/beamer}
		\item[CTAN] \uurl{https://ctan.org/pkg/beamer}
	\end{description}
	\item even though other styles for \LaTeX{} presentations exist, Beamer is by far the best way today
\end{itemize}
%
\end{frame}
%/////////

%/////////
\begin{frame}[c]{AMS}
%
\begin{description}
	\item[AMS] Alma Mater Studiorum--Università di Bologna
	\item[Bologna FC] Bologna Football Club 1909, the city's team
	\begin{itemize}
		\item hence the colours, the \textcolor{bolognafcred}{\textit{rosso}}\textcolor{bolognafcblue}{\textit{bl\`u}}
	\end{itemize}
	\item[Cesena] Campus of Cesena
	\begin{itemize}
		\item hence, this deck having asked for it, the \emph{Campus di Cesena} seal watermarking every slide
	\end{itemize}
\end{description}
%
\end{frame}
%/////////

%/////////
\begin{frame}[c]{Aim}
%
\begin{itemize}
	\item many occasions for presentation
	\item many different people in the group
	\item a common reference could be useful
	\item hence, this style
	\item in the colours of Bologna FC
	\begin{itemize}
		\item the same style as \alert{AMSUniBo} in everything else---options, commands, bibliography
		\item but the faces, which stay \LaTeX{}'s own Computer Modern
	\end{itemize}
\end{itemize}
%
\end{frame}
%/////////

%/////////
\begin{frame}[c]{Who}
%
\begin{itemize}
	\item mostly, myself and my collaborators
	\item everybody else, be my guest
\end{itemize}
%
\end{frame}
%/////////

%/////////
\begin{frame}[c]{Distribution}
%
\begin{block}{Public on GitHub}
	\begin{itemize}
		\item the style is distributed on \uurl{https://github.com/andreaomicini/beamer-AMSBolognaFC/}
		\item the releases are available at \uurl{https://github.com/andreaomicini/beamer-AMSBolognaFC/releases}
		\item this template is distributed on \uurl{https://github.com/andreaomicini/beamer-AMSBolognaFC-template/}
	\end{itemize}
\end{block}
%
\end{frame}
%/////////

%===============================================================================
\section{Use}
%===============================================================================

%/////////
\begin{frame}[c,fragile]{Declaration}
%
\begin{block}{\texttt{\textbackslash{}documentclass}}
Your main Beamer \texttt{.tex} file should open with the declaration
%
\begin{verbatim}
    \documentclass[presentation]{beamer}
        \mode<presentation>{\usetheme{AMSBolognaFC}}
\end{verbatim}
%
so as to use the AMSBolognaFC Beamer style
\end{block}
%
\end{frame}
%/////////

%/////////
\begin{frame}[c,fragile]{Bibliography Style}
%
\begin{block}{\texttt{apalike-AMS}}
Your main Beamer \texttt{.tex} file should include the declaration
\begin{verbatim}
    \bibliographystyle{apalike-AMS}
\end{verbatim}
so as to use the AMSBolognaFC \BibTeX{} style
\end{block}
%
\end{frame}
%/////////

%/////////
\begin{frame}[c,fragile]{Template}
%
\begin{block}{\texttt{AMSBolognaFC-template.tex}}
This template's sources can be used as a simple example of how to use this Beamer style
\end{block}
%
\end{frame}
%/////////


%===============================================================================
\section{Style}
%===============================================================================

%/////////
\begin{frame}[c,fragile]{Colours}
%
The palette is Bologna FC's, and every colour has its own role:
\textcolor{bolognafcblue}{\textit{blue}} structures, \textcolor{bolognafcred}{\textit{red}} marks,
\textcolor{bolognafcsilver!70!black}{\textit{silver}} reads over the dark surfaces
%
\begin{exampleblock}{The four that carry the style}
\footnotesize
\begin{tabular}{@{}llll@{}}
\swatch{bolognafcblue}   & \texttt{bolognafcblue}   & \texttt{\#1A2F48} & text, title, block headers, emphasis, \textsc{toc} \\
\swatch{bolognafcred}    & \texttt{bolognafcred}    & \texttt{\#A21C26} & author, bullets, alerts, links \\
\swatch{bolognafcsilver} & \texttt{bolognafcsilver} & \texttt{\#ECECEC} & over dark: header titles, \texttt{\textbackslash{}ccite}, \texttt{\textbackslash{}aalert} \\
\swatch{bolognafcwhite}  & \texttt{bolognafcwhite}  & \texttt{\#FFFFFF} & the background \\
\end{tabular}
\end{exampleblock}
%
\begin{itemize}
	\item all four are defined by \texttt{beamercolorthemebolognafc.sty}, which the theme loads
	\begin{itemize}
		\item so a deck can use them as any other colour, as in \verb|\textcolor{bolognafcred}{...}|
	\end{itemize}
\end{itemize}
%
\end{frame}
%/////////

%/////////
\begin{frame}[c,fragile]{Blocks}
%
\begin{block}{This is a \texttt{block} environment}
\begin{verbatim}
\begin{block}{...}
...
\end{block}
\end{verbatim}
\end{block}
%
\begin{exampleblock}{This is an \texttt{exampleblock} environment}
\begin{verbatim}
\begin{exampleblock}{...}
...
\end{exampleblock}
\end{verbatim}
\end{exampleblock}
%
\begin{alertblock}{This is an \texttt{alertblock} environment}
\begin{verbatim}
\begin{alertblock}{...}
...
\end{alertblock}
\end{verbatim}
\end{alertblock}
%
\end{frame}
%/////////

%/////////
\begin{frame}[c]{Surfaces}
%
\begin{itemize}
	\item the headers of \texttt{block} and \texttt{alertblock} are \emph{dark} surfaces
	\item everything else is \emph{light}\textemdash{}running text, block bodies, the \texttt{exampleblock} header
	\item since style 1.11, the commands of this style take their colour from the surface they land on
	\begin{itemize}
		\item so a deck never picks a command for its background
	\end{itemize}
	\item a doubled command is the same command, \alert{smaller and quieter}
	\begin{itemize}
		\item \texttt{\textbackslash{}uurl} and \texttt{\textbackslash{}uuurl}, \texttt{\textbackslash{}ddoi} and \texttt{\textbackslash{}dddoi}, \texttt{\textbackslash{}apice} and \texttt{\textbackslash{}aapice}, \ldots
	\end{itemize}
	\item \texttt{\textbackslash{}emph} and \texttt{\textbackslash{}alert} read on every surface, block headers included
\end{itemize}
%
\end{frame}
%/////////

%/////////
\begin{frame}[c,fragile]{One Command, Every Surface}
%
The same source line, here and in each header and body below:\\
\emph{emph}, \alert{alert}, \aalert{aalert}\ccite{bibtex-patashnik88}, \uurl{apice.unibo.it}, \uuurl{apice.unibo.it}
%
\begin{block}{\emph{emph}, \alert{alert}, \aalert{aalert}\ccite{bibtex-patashnik88}, \uurl{apice.unibo.it}, \uuurl{apice.unibo.it}}
\emph{emph}, \alert{alert}, \aalert{aalert}\ccite{bibtex-patashnik88}, \uurl{apice.unibo.it}, \uuurl{apice.unibo.it}
\end{block}
%
\begin{alertblock}{\emph{emph}, \alert{alert}, \aalert{aalert}\ccite{bibtex-patashnik88}, \uurl{apice.unibo.it}, \uuurl{apice.unibo.it}}
\emph{emph}, \alert{alert}, \aalert{aalert}\ccite{bibtex-patashnik88}, \uurl{apice.unibo.it}, \uuurl{apice.unibo.it}
\end{alertblock}
%
\begin{exampleblock}{\emph{emph}, \alert{alert}, \aalert{aalert}\ccite{bibtex-patashnik88}, \uurl{apice.unibo.it}, \uuurl{apice.unibo.it}}
\emph{emph}, \alert{alert}, \aalert{aalert}\ccite{bibtex-patashnik88}, \uurl{apice.unibo.it}, \uuurl{apice.unibo.it}
\end{exampleblock}
%
{\tiny\begin{verbatim}
\emph{emph}, \alert{alert}, \aalert{aalert}\ccite{bibtex-patashnik88},
\uurl{apice.unibo.it}, \uuurl{apice.unibo.it}
\end{verbatim}} %chktex 31
%
\end{frame}
%/////////

%===============================================================================
\section{New Commands}
%===============================================================================

%/////////
\begin{frame}[c,allowframebreaks,fragile]{Citations}
%
\begin{block}{\texttt{\textbackslash{}ccite} command\textemdash{}e.g., \ccite{bibtex-patashnik88}}
\begin{verbatim}
\ccite{bibtex-patashnik88}
\end{verbatim}
\begin{itemize}
	\item to be used instead of standard \texttt{\textbackslash{}cite} command
	\item its colour comes from the surface: light on a dark header, as here above
	\item can be used as a note\ccite{bibtex-patashnik88}, with no space before\textemdash{}faint in running text, by design
	\item since style 1.8 it takes the optional note of \texttt{\textbackslash{}cite}\ccite[p.~42]{bibtex-patashnik88}
	\begin{itemize}
		\item as here above, from \verb|\ccite[p.~42]{bibtex-patashnik88}|
	\end{itemize}
\end{itemize}
\end{block}
%
\begin{exampleblock}{\texttt{\textbackslash{}ccite} in an \texttt{exampleblock} header\textemdash{}e.g., \ccite[p.~42]{bibtex-patashnik88}}
\begin{itemize}
	\item darker on the tinted header, so that it still reads
\end{itemize}
\end{exampleblock}
%
\begin{alertblock}{\texttt{\textbackslash{}cccite} command\textemdash{}e.g., \cccite{bibtex-patashnik88}}
\begin{itemize}
	\item since style 1.11 the same as \texttt{\textbackslash{}ccite}
	\begin{itemize}
		\item a citation marker is already as small as it can be read
		\item kept so that existing decks build as before
	\end{itemize}
\end{itemize}
\end{alertblock}
%
\end{frame}
%/////////

%/////////
\begin{frame}[c,fragile]{URLs}
%
\begin{block}{\texttt{\textbackslash{}uurl} command\textemdash{}e.g., \uurl{https://apice.unibo.it}}
\begin{itemize}
	\item to be used instead of standard \texttt{\textbackslash{}url} command
	\item[e.g.] \verb|\uurl{https://apice.unibo.it}| prints as \uurl{https://apice.unibo.it}
	\item[] and as in this header, light on a dark surface
	\item since style 1.7.5 no character in the address needs escaping
\end{itemize}
\end{block}
%
\begin{block}{\texttt{\textbackslash{}uuurl} command\textemdash{}e.g., \uuurl{https://apice.unibo.it}}
\begin{itemize}
	\item the same, smaller and quieter
	\item[e.g.] \verb|\uuurl{https://apice.unibo.it}| prints as \uuurl{https://apice.unibo.it}
	\item[] and as in this header
\end{itemize}
\end{block}
%
\end{frame}
%/////////

%/////////
\begin{frame}[c,fragile]{DOI}
%
\begin{block}{\texttt{\textbackslash{}ddoi} command\textemdash{}e.g., \ddoi{10.1007/978-3-032-22940-3}} % chktex 8
\begin{itemize}
	\item to be used for the styled, linked DOIs
	\item[e.g.] \verb|\ddoi{10.1007/978-3-032-22940-3}| prints as \ddoi{10.1007/978-3-032-22940-3} % chktex 8
	\item[e.g.] since style 1.7.5 an underscore needs no escaping: \verb|\ddoi{10.1007/978-3-032-22940-3_12}| prints as \ddoi{10.1007/978-3-032-22940-3_12} % chktex 8
\end{itemize}
\end{block}
%
\begin{block}{\texttt{\textbackslash{}dddoi} command\textemdash{}e.g., \dddoi{10.1007/978-3-032-22940-3}} % chktex 8
\begin{itemize}
	\item the same, smaller and quieter\textemdash{}as in the references
	\item[e.g.] \verb|\dddoi{10.1007/978-3-032-22940-3}| prints as \dddoi{10.1007/978-3-032-22940-3} % chktex 8
\end{itemize}
\end{block}
%
\end{frame}
%/////////

%/////////
\begin{frame}[c,fragile]{APICe Addresses}
%
\begin{block}{\texttt{\textbackslash{}apice} command\textemdash{}an object, by its kind and its page name}
\begin{itemize}
	\item[e.g.] \verb|\apice{pub}{AcltFrocos1996}| prints as \apice{pub}{AcltFrocos1996}
	\item[e.g.] \verb|\apice[the AGoT talk]{talk}{AgotAixia2026}| prints as \apice[the AGoT talk]{talk}{AgotAixia2026}
\end{itemize}
\end{block}
%
\begin{block}{\texttt{\textbackslash{}apicesite} and \texttt{\textbackslash{}apicepage} commands}
\begin{itemize}
	\item[e.g.] \verb|\apicesite{andreaomicini}| prints as \apicesite{andreaomicini}
	\item[e.g.] \verb|\apicepage{Courses/Series/Iag}| prints as \apicepage{Courses/Series/Iag}
\end{itemize}
\end{block}
%
\begin{block}{\texttt{\textbackslash{}aapice} and the other doubled forms\textemdash{}e.g., \aapice{course}{Iag2627}}
\begin{itemize}
	\item the same, smaller and quieter, as \verb|\aapice{course}{Iag2627}| in this header; all six since style 1.9
\end{itemize}
\end{block}
%
\end{frame}
%/////////

%/////////
\begin{frame}[c,fragile]{Alert \& Emphasis}
%
\begin{block}{\texttt{\textbackslash{}alert} and \texttt{\textbackslash{}emph} in a header\textemdash{}e.g., \alert{alerted} and \emph{emphasised} text}
\begin{itemize}
	\item both read on every surface: \texttt{\textbackslash{}alert} since style 1.11, \texttt{\textbackslash{}emph} since 1.10.1
	\begin{itemize}
		\item in a header, \texttt{\textbackslash{}emph} is italic in the header's own colour
	\end{itemize}
	\item[e.g.] \verb|\alert{alerted} and \emph{emphasised} text| prints as \alert{alerted} and \emph{emphasised} text
\end{itemize}
\end{block}
%
\begin{block}{\texttt{\textbackslash{}aalert} command\textemdash{}e.g., \aalert{quietly alerted text}}
\begin{itemize}
	\item a quiet alert, italic, on every surface
	\item[e.g.] \verb|\aalert{quietly alerted text}| prints as \aalert{quietly alerted text}
\end{itemize}
\end{block}
%
\end{frame}
%/////////

%/////////
\begin{frame}[c,fragile,allowframebreaks]{Speakers \emph{vs.} Authors}
%
\begin{exampleblock}{\texttt{\textbackslash{}speaker} command---e.g., \speaker{Diego Zorro}}
\begin{itemize}
	\item to be used within \texttt{\textbackslash{}author} standard Beamer command to single out the actual speaker among the authors
	\item in bold, and in a colour of its own
	\item[e.g.] as in
\begin{verbatim}
	\author[Garcia \and Zorro]
	{Sarg Garcia \and \speaker{Diego Zorro}}
\end{verbatim}
	\item and in the author specification of this template
\end{itemize}
\end{exampleblock}
%
\begin{block}{\texttt{\textbackslash{}sspeaker} command---e.g., \sspeaker{Diego Zorro}}
\begin{itemize}
	\item to be used within \texttt{\textbackslash{}author} standard Beamer command to single out the actual speaker among the authors in the short form
	\item in bold, and in the colour of the surface it sits on---light on the footline, as below
	\item[e.g.] as in
\begin{verbatim}
	\author[Garcia \and \sspeaker{Zorro}]
	{Sarg Garcia \and \speaker{Diego Zorro}}
\end{verbatim}
	\item and in the author specification of this template
\end{itemize}
\end{block}

\end{frame}
%/////////

%===============================================================================
\section{New Options}
%===============================================================================

%/////////
\begin{frame}[c,fragile,allowframebreaks]{\texttt{apice} option}
%
\begin{alertblock}{\texttt{apice} Beamer option}
\begin{itemize}
	\item the option \texttt{apice} can be used in the presentation preamble
	\begin{itemize}
		\item as in
		\begin{verbatim}
\documentclass[presentation,apice]{beamer}...
		\end{verbatim}
	\end{itemize}
	\item enables the display of the \BibTeX{} field \texttt{apice}
	\begin{itemize}
		\item containing the relative URL of the publication on the \alert{\href{https://apice.unibo.it/}{\textsf{APICe}}} XWiki
	\end{itemize}
	\item properly linked to the corresponding \href{https://apice.unibo.it/}{\textsf{APICe}} page
\end{itemize}
\end{alertblock}
%
\begin{exampleblock}{Example}
%
The following \BibTeX{} entry in the \texttt{.bib} file with the \texttt{apice} option enabled
%
{\tiny\begin{verbatim}
@manual{bibtex-patashnik88,
    apice = {BibtexPatashnik88},
    author = {Patashnik, Oren},
    booktitle = {\BibTeXing},
    organization = {CTAN, The Comprehensive TeX Archive Network},
    title = {\BibTeX{}ing},
    url = {https://mirrors.ctan.org/biblio/bibtex/base/btxdoc.pdf},
    year = 1988}
\end{verbatim}} %chktex 31
%
would translate into the following entry in the presentation bibliography
%
\begin{thebibliography}{}
\bibitem[Patashnik, 1988]{bibtex-patashnik1988}
Patashnik, O. (1988)
\newblock {\bfseries {\em \BibTeX{}ing}}
\newblock CTAN, The Comprehensive TeX Archive Network
\newblock\apicepubpar{BibtexPatashnik88} \uuurl{https://mirrors.ctan.org/biblio/bibtex/base/btxdoc.pdf}
\end{thebibliography}
%
\end{exampleblock}
%
\end{frame}
%/////////

%/////////
\begin{frame}[c,fragile]{\texttt{bologna}, \texttt{cesena} options}
%
\begin{alertblock}{Which seal watermarks the slides---none, unless asked for}
\begin{itemize}
	\item \texttt{bologna}---the seal of the \emph{Alma Mater Studiorum} as a whole
	\item \texttt{cesena}---the seal of the \emph{Campus di Cesena}
	\begin{itemize}
		\item the two are told apart by the border, solid for the Campus and an outline for the Ateneo
	\end{itemize}
	\item both sit in the lower-right corner, cropped by the page edge
	\begin{itemize}
		\item the one on this very slide is the one this deck selected
	\end{itemize}
	\item \alert{neither option, no seal}---the slides stay bare
	\begin{itemize}
		\item the seals are institutional marks, so a deck claims one out loud
		\item which leaves the style free to use for whoever does not belong---or no longer, or not yet belongs---to the \emph{Alma Mater}
	\end{itemize}
\end{itemize}
\end{alertblock}
%
\begin{exampleblock}{Example}
{\footnotesize\begin{verbatim}
\documentclass[presentation,cesena]{beamer}
    \mode<presentation>{\usetheme{AMSBolognaFC}}
\end{verbatim}} %chktex 31
\end{exampleblock}
%
\end{frame}
%/////////

%===============================================================================
\section*{}
%===============================================================================

%/////////
\frame{\titlepage}
%/////////

%===============================================================================
\section*{\refname}
%===============================================================================

%%%%
\setbeamertemplate{page number in head/foot}{}
%/////////
\begin{frame}[c,noframenumbering]{\refname}
%\begin{frame}[t,allowframebreaks,noframenumbering]{\refname}
%	\tiny
	\scriptsize
%	\footnotesize
	\bibliographystyle{apalike-AMS}
	\bibliography{AMSBolognaFC-template}
\end{frame}
%/////////

%%%%%%%%%%%%%%%%%%%%%%%%%%%%%%%%%%%%%%%%%%%%%%%%%%%%%%%%%%%%%%%%%%%%%%%%%%%%%%%%
\end{document}
%%%%%%%%%%%%%%%%%%%%%%%%%%%%%%%%%%%%%%%%%%%%%%%%%%%%%%%%%%%%%%%%%%%%%%%%%%%%%%%%
